\section{Definície}
V tejto časti sa inšpirujeme relačnou algebrou z \cite{RAcompilator} a klasickou \textit{učebnicovou} relačnou algebrou a definujeme našu vlastnú.

\subsection{Učebnicová relačná algebra}
Klasická relačná algebra pracuje s reláciami, ktoré majú väčšinou jednotlivé atribúty pomenované. Potom ich vieme jednoducho reprezentovať pomocou tabuliek.

\begin{priklad}
    \[
    \begin{array}{c c}
        \begin{aligned}
            R &\in \textbf{Meno} \times \textbf{Výška} \\
            R &= \{(\text{Adam}, 175), (\text{Boris}, 183), (\text{Cecil}, 192)\}
        \end{aligned}
        &
        \begin{array}{c|c}
            \multicolumn{2}{c}{R} \\ \hline
            \textbf{Meno}  & \textbf{Výška} \\ \hline
            \text{Adam}  & 175 \\
            \text{Boris} & 183 \\
            \text{Cecil} & 192
        \end{array}
    \end{array}
    \]
\end{priklad}

Na manipuláciu s týmito reláciami existuje niekoľko operátorov:
\begin{itemize}
    \item \textbf{Množinové operátory} \\ 
    zjednotenie množín $\cup$, rozdiel množín $-$, (prienik množín $\cap$, ale ten sa dá vyjadriť aj za pomoci rodielu $A \cap B = A - (A-B)$)
    
    Pre relácie, na ktoré aplikujeme tieto operátory musí platiť, že majú identické domény atribútov. Pred tým, ako aplikujeme operátor na dve relácie, musíme usporiadať ich atribúty v jednotlivých záznamoch do rovnakého poradia.
    
    \item \textbf{Karteziánsky súčin $\times$} \\
    $R \times S$, kde $R$ a $S$ sú relácie. Výsledná relácia sa dá matematicky zapísať nasledovne: 
    $$R \times S = \{(a_1, \dots, a_n, b_1, \dots, b_m) \, | \,  (a_1, \dots, a_n) \in R, (b_1, \dots, b_m) \in S)\}$$
    \begin{priklad}
    \begin{center}
        \[
        \begin{array}{c c c}
            \begin{array}{c|c}
            \multicolumn{2}{c}{R} \\ \hline
            \textbf{X} & \textbf{Y} \\ \hline
            A & a\\
            B & b \\
            C & c
            \end{array}
            &
            \begin{array}{c|c}
            \multicolumn{2}{c}{S} \\ \hline
            \textbf{Y} & \textbf{Z} \\ \hline
            1 & 10\\
            2 & 20\\
            3 & 30
            \end{array}
            &
            \begin{array}{c|c|c|c}
            \multicolumn{4}{c}{R\times S} \\ \hline
            \textbf{X} & \textbf{R.Y} & \textbf{S.Y} & \textbf{Z} \\ \hline
            A & a & 1 & 10\\
            B & b & 1 & 10\\
            C & c & 1 & 10\\
            A & a & 2 & 20\\
            B & b & 2 & 20\\
            C & c & 2 & 20\\
            A & a & 3 & 30\\
            B & b & 3 & 30\\
            C & c & 3 & 30
            \end{array}
        \end{array}
        \]
    \end{center}
    \end{priklad}
    
    \item \textbf{Projekcia $\Pi$} \\
    $\Pi_{a_1, \dots, a_n}(R)$, kde $a_1, \dots, a_n$ sú názvy atribútov relácie $R$ vráti novú reláciu, ktorá obsahuje \uv{iba stĺpce} $a_1, \dots, a_n$. Výsledok je stále relácie (čiže aj množina), takže neobsahuje duplicitné záznamy.
    \begin{priklad}
    \begin{center}
        \[
        \begin{array}{c c c}
            \begin{array}{c|c|c}
            \multicolumn{3}{c}{R} \\ \hline
            \textbf{Meno} & \textbf{Mesto}    & \textbf{Zviera} \\ \hline
            \text{Adam}   & \text{Bratislava} & \text{Pes} \\
            \text{Boris}  & \text{Bratislava} & \text{Mačka} \\
            \text{Cecil}  & \text{Košice}     & \text{Pes} \\
            \text{Dano}   & \text{Lučenec}    & \text{Zajac} \\
            \text{Erik}   & \text{Poprad}     & \text{Kôň} \\
            \text{Fero}   & \text{Košice}     & \text{Mačka} \\
            \end{array}
            &
            \begin{array}{c|c}
            \multicolumn{2}{c}{\Pi_{\textbf{Meno}, \textbf{Zviera}}(R)} \\ \hline
            \textbf{Meno} & \textbf{Zviera} \\ \hline
            \text{Adam}   & \text{Pes} \\
            \text{Boris}  & \text{Mačka} \\
            \text{Cecil}  & \text{Pes} \\
            \text{Dano}   & \text{Zajac} \\
            \text{Erik}   & \text{Kôň} \\
            \text{Fero}   & \text{Mačka} \\
            \end{array}
            &
            \begin{array}{c}
            \multicolumn{1}{c}{\Pi_{\textbf{Mesto}}(R)} \\ \hline
            \textbf{Mesto}   \\ \hline
            \text{Bratislava}\\
            \text{Košice}    \\
            \text{Lučenec}   \\
            \text{Poprad}    \\
            \end{array}
        \end{array}
        \]
    \end{center}
    \end{priklad}
    
    \item \textbf{Selekcia $\sigma$} \\
    $\sigma_C(R)$, kde $C$ je podmienka, ktorá obsahuje atribúty relácie $R$. Výsledná relácie je podmnožina relácie $R$ obsahujúca iba tie záznamy, ktoré splňujú podmienku $C$.
    \begin{priklad}
    \begin{center}
        \[
        \begin{array}{c c}
            \begin{array}{c|c|c}
            \multicolumn{3}{c}{R} \\ \hline
            \textbf{Meno} & \textbf{Mesto}    & \textbf{Zviera} \\ \hline
            \text{Adam}   & \text{Bratislava} & \text{Pes} \\
            \text{Boris}  & \text{Bratislava} & \text{Mačka} \\
            \text{Cecil}  & \text{Košice}     & \text{Pes} \\
            \text{Dano}   & \text{Lučenec}    & \text{Zajac} \\
            \text{Erik}   & \text{Poprad}     & \text{Kôň} \\
            \text{Fero}   & \text{Košice}     & \text{Mačka} \\
            \end{array}
            &
            \begin{array}{c|c|c}
            \multicolumn{3}{c}{\sigma_{\textbf{Mesto} = \text{Bratislava} \lor \textbf{Zviera} = \text{Mačka}}(R)} \\ \hline
            \textbf{Meno} & \textbf{Mesto}    & \textbf{Zviera} \\ \hline
            \text{Adam}   & \text{Bratislava} & \text{Pes} \\
            \text{Boris}  & \text{Bratislava} & \text{Mačka} \\
            \text{Fero}   & \text{Košice}     & \text{Mačka} \\
            \end{array}
        \end{array}
        \]
    \end{center}
    \end{priklad}
    
    \item \textbf{Prirodzený (natural) join $\Join$} \\
    $R \Join S$, kde $R$ a $S$ sú relácie. Podobne ako v karteziánskom súčine, párujeme záznamy z dvoch relácíí, ale párujeme iba tie záznamy z $R$ a $S$, ktoré sa zhodujú na tých atribútoch, ktoré sú spoločné v $R$ a $S$. V množinovom zápise:
    \begin{align*}
        R          &= A_1 \times \cdots \times A_n \times C_1 \times \cdots \times C_k\\
        S          &= B_1 \times \cdots \times B_n \times C_1 \times \cdots \times C_k\\
        R \times S &= \{(a_1, \dots, a_n, c_1, \dots, c_k, b_1, \dots, b_m) \, | \, &&(a_1, \dots, a_n, c_1, \dots, c_k) \in R, \\
        &                                                                &&(b_1, \dots, b_m, c'_1, \dots, c'_k) \in S, \\
        &                                                                &&\ c_1 = c'_1, \dots, c_k = c'_k\}
    \end{align*}
    \begin{priklad}
        \begin{center}
            TODO
        \end{center}
    \end{priklad}
    
    \item \textbf{Antijoin $\Antijoin$} \\
    $R \Antijoin S$, kde $R$ a $S$ sú relácie. Vo výslednej relácií sú tie záznami z $R$, pre ktoré neexistuje záznam v $S$, s ktorým by sa zhodovali na spoločných atribútoch. V množinovom zápise: 
    \begin{align*}
        R          &= A_1 \times \cdots \times A_n \times C_1 \times \cdots \times C_k\\
        S          &= B_1 \times \cdots \times B_n \times C_1 \times \cdots \times C_k\\
        R \times S &= \{(a_1, \dots, a_n, c_1, \dots, c_k) \in R \, | \, \forall (b_1, \dots, b_m, c'_1, \dots, c'_k) \in S: c_1 \neq c'_1, \dots, c_k \neq c'_k\}
    \end{align*}
    \begin{priklad}
        \begin{center}
            TODO
        \end{center}
    \end{priklad}
    
    \item \textbf{Theta-join $\Join_C$} \\
    $R \Join_C S$, kde $R$ a $S$ sú relácie a $C$ je podmienka, v ktorej sa vyskytujú atribúty z $R$ a $S$. Natural join nám povolil spojiť dva záznamy iba ak ich spoločné atribúty boli rovnaké. Theta-join popáruje každú dvoji záznamov z $R$ a $S$, ktoré spĺňajú podmienku $C$.
    \begin{priklad}
    \begin{center}
        TODO
    \end{center}
    \end{priklad}
    
    \item \textbf{Premenovanie $\rho$}\\
    $\rho_{A/ B} (R)$, kde $R$ je relácie, $A$ je názov atribútu z relácie $R$ a $B$ je nový názov atribútu. Výsledkom bude relácia s rovnakými záznamami, iba atribút $A$ a premenuje na $B$.
    \begin{priklad}
    \begin{center}
        \[
        \begin{array}{c c}
            \begin{array}{c|c|c}
            \multicolumn{3}{c}{R} \\ \hline
            \textbf{Meno} & \textbf{Mesto}    & \textbf{Zviera} \\ \hline
            \text{Adam}   & \text{Bratislava} & \text{Pes} \\
            \text{Boris}  & \text{Bratislava} & \text{Mačka} \\
            \text{Cecil}  & \text{Košice}     & \text{Pes} \\
            \text{Dano}   & \text{Lučenec}    & \text{Zajac} \\
            \text{Erik}   & \text{Poprad}     & \text{Kôň} \\
            \text{Fero}   & \text{Košice}     & \text{Mačka} \\
            \end{array}
            &
            \begin{array}{c|c|c}
            \multicolumn{3}{c}{\rho_{\textbf{Mesto} / \textbf{Obec}}(R)} \\ \hline
            \textbf{Meno} & \textbf{Obec}     & \textbf{Zviera} \\ \hline
            \text{Adam}   & \text{Bratislava} & \text{Pes} \\
            \text{Boris}  & \text{Bratislava} & \text{Mačka} \\
            \text{Cecil}  & \text{Košice}     & \text{Pes} \\
            \text{Dano}   & \text{Lučenec}    & \text{Zajac} \\
            \text{Erik}   & \text{Poprad}     & \text{Kôň} \\
            \text{Fero}   & \text{Košice}     & \text{Mačka} \\
            \end{array}
        \end{array}
        \]
    \end{center}
    \end{priklad}
    
    \item \textbf{Iné join-y} \\
    V niektorých knihách sa spomínajú aj iné variácie join-ov ako \textit{left outer join}, \textit{right outer join} alebo \textit{full outer join}. My sa s nimi nebudeme zaoberať, keďže sa dajú vyskladať z predchádzajúcich operátorov.
    Pre ich správne fungovanie treba ale zaviesť špeciálnu konštantu, ktorá sa v literatúre označuje ako $\omega$, $\varepsilon$, \verb|null| alebo \verb|none|.
\end{itemize}

\subsection{Naša relačná algebra}
Pre jednoduchší preklad datalogu do relačnej algebry sme si zvolili nasledujúcu algebru s piatimi operátormi (inšpirovaná algebrou z \cite{RAcompilator}). Tiež pracuje s reláciami, ale nepotrebuje vopred pomenované atribúty. Vždy, keď potrebuje pracovať s \uv{názvamy stĺpcov}, sú explicitne vymenované.

\begin{itemize}
    \item
    {
        $\UNION([R_1, \dots, R_n])$, kde $R_1, \dots, R_n$ sú relácie. Výsledná relácia bude obsahovať všetky záznamy z každej relácie $R_1$ až $R_n$.
    
        \begin{priklad}
        \begin{center}
            \[
        \begin{array}{c c c c}
            \begin{array}{c|c}
            \multicolumn{2}{c}{A} \\ \hline
            11 & a \\
            12 & b \\
            13 & c
            \end{array}
            &
            \begin{array}{c|c}
            \multicolumn{2}{c}{B} \\ \hline
            21 & a \\
            22 & b \\
            23 & c
            \end{array}
            &
            \begin{array}{c|c}
            \multicolumn{2}{c}{C} \\ \hline
            31 & a \\
            32 & b \\
            33 & c
            \end{array}
            &
            \begin{array}{c|c}
            \multicolumn{2}{c}{\UNION([A, B, C])} \\ \hline
            11 & a \\
            12 & b \\
            13 & c \\
            21 & a \\
            22 & b \\
            23 & c \\
            31 & a \\
            32 & b \\
            33 & c 
            \end{array}
        \end{array}
        \]
        \end{center}
        \end{priklad}
    } 
    
    \item
    {
        $\JOIN([R_1, \dots, R_n], [[T_{1, 1}, \dots, T_{1, a_1}], \dots, [T_{n, 1}, T_{n, a_n}]], [T_1, \dots, T_k])$,
         kde $R_i$ je relácia s aritou $a_i$, $T_{i, j}$ sú \textit{vstupné} termy a $T_l$ sú \textit{výstupné} termy.
        
        Výsledná relácia bude nasledujúca. Nech $P_1, \dots, P_p$ sú všetky premenné vyskytujúce sa v termoch $T_{1, 1}, \dots, T_{1, a_1}, T_{2, 1}, \dots, T_{n, a_n}$. Nech mnoňina $M$ obsahuje všetky $(c_1, \dots, c_p)$ také, pre ktoré platí 
        $$(T_{i, 1}[P_1|c_1, \dots, P_p|c_p], \dots, T_{i, a_i}[P_1|c_1, \dots, P_p|c_p]) \in R_i, \, \forall i \in \{1, \dots, n\}\text{.}$$
        Potom vo výslednej relácií $R$ budú záznamy 
        $$R = \{(T_1[P_1|k_1, \dots, P_p|k_p], \dots, T_k[P_1|k_1, \dots, P_p|K_p]) \, | \, (k_1, \dots, k_p) \in M\} \text{.}$$
        
        \begin{priklad}
        \begin{center}
            \[
            \begin{array}{c c}
                \begin{array}{c|c}
                    \multicolumn{2}{c}{G} \\ \hline
                    1 & 2 \\
                    2 & 3 \\
                    3 & 4 \\
                    3 & 5 \\
                \end{array}
                &
                \begin{array}{c|c}
                    \multicolumn{2}{c}{\JOIN([G, G], [[X, Z], [Z, Y]], [X, Y])} \\ \hline
                    1 & 3 \\
                    2 & 4 \\
                    2 & 5 \\
                \end{array}
            \end{array}
            \]
        \end{center}
        \end{priklad}
    }
        
    \item
    {
        $\ANTIJOIN([R_1, \dots, R_n], [[T_{1, 1}, \dots, T_{1, a_1}], \dots, [T_{n, 1}, T_{n, a_n}]], [T_1, \dots, T_k])$,
        kde $R_i$ je relácia s aritou $a_i$, $T_{i, j}$ sú \textit{vstupné} termy a $T_l$ sú \textit{výstupné} termy.
        
        Výsledná relácia bude nasledujúca. Nech $P_1, \dots, P_p$ sú všetky premenné vyskytujúce sa v termoch $T_{1, 1}, \dots, T_{1, a_1}, T_{2, 1}, \dots, T_{n, a_n}$. Nech mnoňina $M$ obsahuje všetky $(c_1, \dots, c_p)$ také, pre ktoré platí 
        $$(T_{1, 1}[P_1|c_1, \dots, P_p|c_p], \dots, T_{1, a_1}[P_1|c_1, \dots, P_p|c_p]) \in R_1$$
        a zároveň pre aspoň jedno $i \in \{2, \dots, n\}$ platí
        $$(T_{i, 1}[P_1|c_1, \dots, P_p|c_p], \dots, T_{i, a_i}[P_1|c_1, \dots, P_p|c_p]) \notin R_i\text{.}$$
        Potom vo výslednej relácií $R$ budú záznamy 
        $$R = \{(T_1[P_1|k_1, \dots, P_p|k_p], \dots, T_k[P_1|k_1, \dots, P_p|K_p]) \, | \, (k_1, \dots, k_p) \in M\} \text{.}$$
        \begin{priklad}
        \begin{center}
            TODO
        \end{center}
        \end{priklad}
    }
        
    \item 
    {
        $\REC(N, R)$, kde $N$ je názov relácie a $R$ je relácia. Tento operátor definuje reláciu s názvom $N$ a dovoľuje použiť v relácií $R$ odkaz na ňu.
        \begin{priklad}
        \begin{center}
            $$\REC
            (
                t, 
                \UNION
                (
                    [
                        \$s,
                        \JOIN
                        (
                            [t, t],
                            [[X, Z], [Z, Y]],
                            [X, Y]
                        )
                    ]
                )
            )$$
            Relácia $t$ je tranzitívny uzáver vstavanej relácie $s$.
        \end{center}
        \end{priklad}
    }
    
    \item
    {
        $\DEF([[T_{1, 1}, \dots, T_{1, n}], \dots, [T_{k, 1}, \dots, T_{k, n}]])$, kde $T_{1, 1}, \dots, T_{k, n}$ sú termy \textbf{bez} premenných. Výsledok bude relácia $R = \{(T_{i, 1}, \dots, T_{i, n}) \, | \, i \in \{1, \dots, k\}\}$
        \begin{priklad}
        \begin{center}
            TODO
        \end{center}
        \end{priklad}
    }
\end{itemize}
